\documentclass{slides}

\title{\includegraphics[width=0.5\paperwidth]{graphics/war.pdf}}
\subtitle{WAR with Auction Rounds}
\author{Jack Rosenthal, Paul Sattizahn}
\institute{CSCI-410: Game Theory \& Networks}

\let\leq=\leqslant
\let\geq=\geqslant

\usepackage{pgfplots}

\begin{document}

\begin{frame}{Assumptions}
    Since:
    \begin{enumerate}
        \item The starting state of each player is known to all other players
        \item Every player knows what all other players played
    \end{enumerate}

    \bigskip

    We may assume:

    \begin{block}{Perfect Information Principle}
        Since every player is capable of ``counting cards'', every player is
        theoretically capable of knowing the cards in every other player $i$'s
        hand ($H_i$) and wins pile ($W_i$).
    \end{block}

\end{frame}

\begin{frame}{WAR as a Repeated Static Game}
    WAR can be represented as a repeated static game; that is, the outcome of a
    player's move depends on the other players' moves, but state is maintained
    between moves to formulate the outcome of game.
\end{frame}

\begin{frame}{Problem Statement}
    You already know how WAR works.

    \pause
    \bigskip

    Given a game of WAR with $N$ players, devise a \textbf{strategy} to ``do
    well'' at WAR.
\end{frame}

\begin{frame}{Dead Simple Strategy: DummiePlayer (DP)}

    \textbf{DummiePlayer$i$ (DP)}:
    \begin{itemize}
        \item To play a normal move, select randomly from $H_i$
        \item To play a tie, choose 4 cards randomly from $H_i$ and discard the
            first 3, play the last
    \end{itemize}

    \pause
    \bigskip

    \textbf{How well does DP do?} When playing against $N$ other DP's, a
    DP will have a $\frac1N$ chance of winning.

\end{frame}

\begin{frame}{An Intelligent Player: SimpleMindedPlayer (SMP)}
    \textbf{SimpleMindedPlayer$i$ (SMP)}:
    \begin{itemize}
        \item To determine a normal play:
            \begin{enumerate}
                \item Compute $S = \{H_{ic} \,|\, \forall j,d\, [H_{ic} >
                    H_{jd}]\}$.
                \item If $S \ne \emptyset$, play $\min(S)$.
                \item Otherwise if, $|H_i| \geq 4$, $\exists j,c \,[H_{ic} =
                    \max\left(H_{j}\right) \land \max_2\left(H_{i}\right) \geq
                    \max_2\left(H_{j}\right)]$, then play the corresponding $H_{ic}$.
                \item Otherwise, play $\min\left(H_{i}\right)$.
            \end{enumerate}
        \item To play a tie, follow strategy similar to above, but remove 3
              smallest cards from hand first to be discarded.
    \end{itemize}
    \bigskip
    \pause

    \textbf{How well does a SMP do against $N$ other SMPs?} A SMP will always
    die against a fellow SMP, and the game will result with no winners. So an
    SMP stands a 0\% chance here.

    \bigskip
    \pause

    \textbf{How well does a SMP do against $N$ DPs?} Too hard to figure out
    mathematically, so we wrote a computer simulation to answer this.

\end{frame}

% TODO SMP calculation example

\begin{frame}{Computer Simulation}

    \begin{itemize}
        \item 100,000 simulated games
        \item A single SMP battles $N$ DPs, for all $N \in \{1, \ldots, 100\}$.
    \end{itemize}
\end{frame}

\begin{frame}{Computer Simulation}

    \begin{center}
    \begin{tikzpicture}
    \begin{axis}[
            title=SMP vs. $N$ DPs: Win Ratio,
            xlabel={Number of DPs ($N$)},
            ylabel={SMP Win Ratio},
    ]
    \addplot
    [blue] table {wins.dat};
    \end{axis}
    \end{tikzpicture}
    \end{center}
\end{frame}
\begin{frame}{Computer Simulation}

    \begin{center}
    \begin{tikzpicture}
    \begin{axis}[
            title=SMP vs. $N$ DPs: Mean Rounds Played,
            xlabel={Number of DPs ($N$)},
            ylabel={Mean Rounds Played},
    ]
    \addplot
    [red] table {mrounds.dat};
    \end{axis}
    \end{tikzpicture}
    \end{center}
\end{frame}
\begin{frame}{Computer Simulation}

    \begin{center}
    \begin{tikzpicture}
    \begin{axis}[
            title=SMP vs. $N$ DPs: Variance of Rounds Played,
            xlabel={Number of DPs ($N$)},
            ylabel={$\sigma$ (Rounds)},
    ]
    \addplot
    [green] table {vrounds.dat};
    \end{axis}
    \end{tikzpicture}
    \end{center}
\end{frame}

\begin{frame}{Extensions}
    \begin{itemize}[<+->]
        \item SMP does not take into account $W_i$ for either themselves, or
            other players. Design a CMP (ComplexMindedPlayer) which adds in
            factors from $W_i$.
        \item When an SMP determines it cannot win, it chooses it's lowest
            card, which may have a high probability of causing a tie against
            DPs. Design a player good for defeating large quantities of DPs
            that takes this into account.
        \item Throw the game at a ML algorithm and see what happens?
        \item More simulation runs!
    \end{itemize}
\end{frame}

{
    \setbeamercolor{background canvas}{bg=csblue}
\begin{frame}

    \begin{center}
        \Huge
        \color{white}{Questions?}
    \end{center}

\end{frame}
}

\end{document}
